\subsection{Contexte skinsoft}

\begin{itemize}
	
\item L'entreprise 
	
Skinsoft propose des solutions logicielles full web principalement à destination d'institutions patrimoniales. 

L'application générale est divisée en différents espaces de travail, chacun spécifique à un type de collection patriculier.  Par exemple, un musée aura un module S-museum, un musée d'histoire naturelle aura un module Skin Specimen, une cinémathèque aura un module S-movies. Ces espaces de travail sont eux-mêmes configurables et adaptables en fonction des pratiques métiers de l'institution. Skinsoft peut ainsi paramétrer l'utilisation de certains champs, de certains types de notices sur demande de l'utilisateur. 

En amont de la livraison de l'application chez un client, un mapping est effectué entre les champs d'indexation et de catalogage utilisés par le client dans l'ancienne base de données, et les champs que propose skinsoft dans son application. Même si les champs sont configurables, une adaptation et une mise à niveau doit s'effectuer. C'est d'ailleurs le point le plus complexe à gérer en termes de relations clients. Le client doit en effet accepter d'abandonner certaines pratiques pour pouvoir adopter un nouveau logiciel. Ce processus est fastidieux et demande de nombreuses formations pour accompagner les équipes métiers. 

En un mot, l'idée de l'application proposée par skinsoft est celle de la configurabilité et de l'adaptabilité. C'est ce mot d'ordre qui pousse l'entreprise à repenser son logiciel à chaque nouveau contrat, pour être le plus proche possible des pratiques métiers de chaque institution et des différentes attentes. 

Ce qui caractérise l'application de skinsoft, davantage qu'une simple édition de données, c'est surtout la recherche de données et la flexibilité de la recherche de données. 

\item mise en place de l'IA

%à ajouter ici : info sur le projet SKAI : contexte du projet IA chez skinsoft et ligens directrices du projet

La gestion de collections culturelles est un métier très normé, régit par des normes professionnelles pour favoriser le partage des connaissances. L'IA émergente devra être enstraîné sur ces normes professionnelles. 

Avec l'arrivée de l'IA dans le monde du travail et sa remise en question de la place humaine dans certaines pratiques profesionnelles qui seraient automatisables, une entreprise dédiée à l'amélioration des pratiques métier en contexte patrimonial et culturel comme skinsoft souhaite l'implémenter. 

Cette volonté d'implémenter l'IA s'inscrit dans l'appel à projets "France 2030" qui couvre notamment la maitrise des technologies numériques et le "déploiement de l'innovation". \footcite{zotero-item-636}. Ce projets finance des initiatives à l'intersection de la culture, du numérique et de l'IA.

L'une des caractéristiques majeures de l'application, nous l'avons dit, est sa flexibilité et son adaptabilité. La mise en place de l'IA permettrait de mettre cette adaptabilité en place plus efficacement et à plus large échelle. Automatiser la migration de données, le paramétrage de l'application, la navigation dans l'application permettrait à l'entreprise de satisfaire ses clients plus rapidement et de gérer ainsi plusieurs contrats à la fois. Mais l'IA est-elle attendue côté client ? 

À quels besoins concrets chez les clients skinsoft l'IA répond-elle ?
aussi, on parle d'IA mais qu'est ce que ça implique comme notions et outils concrètement ? (voir IAgenerale.tex)

- améliorer l'expérience client dans l'application
- améliorer l'efficacité de l'application
- augmenter l'efficacité des pratiques métier 
- donner accès plus rapidement et à un plus grand nombre d'archives 
- promouvoir une meilleur organisation interne

En effet, la question du besoin est centrale, pour une utilisation de l'IA éthique, comme le mentionne le ministère de la culture : " il ne s’agit pas d’utiliser l’IA par défaut, mais d’abord d’examiner le besoin à couvrir et de vérifier la capacité des systèmes d’IA à satisfaire ce
besoin dans des conditions optimales ; le cas échéant, de favoriser des alternatives moins consommatrices que l’IA si elles sont suffisantes à répondre au besoin". \footcite{2025a}

L'objectif est double pour une entreprise comme skinsoft : 
La mise ne place de l'IA dans l'application permettrait d'autonomisr les institutions sur leur usage du logiciel et d'être plus productifs dans l'indexation de leurs collections. 
Ainsi, les équipes skinsoft seraient moins sollicitées pour des problèmes liés à l'usage de l'application qui pourraient être réglés par un modèle de langage. 
La mise en place d'une barre de recherche fondée sur le langage naturel découle de ces objetcifs, ainsiq ue la mise en pace d'un chatbot, deux idées déjà amorcées et au coeur des discussions chez Skinsoft.

Les fonctionnalités que skinsoft souhaite mettre en place, liées  la vidéo : composants de l'image avec proposition d'indexation, transcription du son et proposition d'indexation sujets, reconnaissance vocale et proposition d'indexation musique, personnage, ambiance. Important : la mise en plce de l'IA se fait par l'identification de données ouvertes qui peuvent nourir le moteur. 
Le projet est de paramétrer et d'entraîner plusieurs moteurs  : socle de savoir. 
le moteur doit il être entraîné sur des donnes publiques, ou entraîné complémentairement par des données de l'institution sur laquelle il travaille. il faut que les clients acceptent : contrat spécifique. (limites) 
ce que veut skinsoft : création de contenus structurés, recommandation de contenus structurés, assistance à l'usage, force de proposition créative, détection...
 
choix d'appui sur des données produites directement par des institutions culturelles. utiliser exclusivement des sources ouvertes. pour les images : entrainement sur des images disponibles sous licence ouverte : european, POP + textes de lois, méthodologie des gestion des collections, thesaurus, normes de catalogage...Restaurant Caspian, Rue des Entrepreneurs, Paris
Il faut l'accord client et les données d'un seul client ne sont pas suffisantes.
l'indexation est un composant parmi d'autres du projet IA chez skinsoft, avec tout ce qui va être mis en place également. 
Le moteur IA doit jouer le rôle de prescripteur, force de proposition, et de facilitateur. (attention : mots exacts employés par skinsoft). Il revient à l'utilisateur de retenir tout ou une partie de la proposition faite par l'IA. 

Mais d'autres besoins métier sont susceptibles d'être assistés par IA, au delà de la recherche et de questionnements d'usage. En effet, skinsoft souhaiterait mettre en place une automatisation de l'indexation, permettant ainsi d'une part d'alléger la charge de travail des équipes au sein d'une institution patrimoniale, et d'autre part, de donner accès à davantage de documents pour le grand public. 

Faire le lien avec l'indexation des archives audiovisuelles : pas les mêmes défis d'indexation pour une collection audiovisuelles : lié à la spécificité de ces documents "boites noires"

\item la tension privé public
Si skinsoft est une entreprise privée, elle adopte en revanche une posture publique culturelle, celle de ses clients, afin de répndre à leurs besoins. La mise en place de l'IA est pensée pour améliorer les logiques patrimoniales quotidiennes telles que l'indexation, le catalogage, le classement, la gestion de projet... Ces pratiques ne sont pas liées à des objectifs de rentabilité et de productivité concurrentiels comme c'est le cas d'entreprises privées. Ainsi, skinsoft fait face à une tension : l'entreprise s'affiche comme une entreprise innovante en matière de technologies numériques, alors même que ses clients sont cantonnés à des pratiques presqu'ancestrales, celle de la gestion de collections. 

Au sein de la mise en place et de l'arrivée de l'IA dans un logiciel de gestion des collections, il faut ainsi se questionner sur sa pertinence. Ces nouvelles fontionnalités sont elles attendues ? Les institutions sont-elles prêtes à utiliser l'IA au quotidien pour gérer leur collection? pour les décrire ? pour les rendre accessible au public ?

En 2025, le ministère de la culture identifie "l'enrichissement de la connaissances et des savoirs" comme un usage potentiel de l'IA en secteur culturel avec par exemple "de l’aide au catalogage, à l’indexation automatique des collections, au tri et au classement des
documents et données destinés à être archivés, à la transcription des textes manuscrits, à la reconnaissance de formes, à la diffusion et à la valorisation des contenus culturels et patrimoniaux quel que soit leur
format (image, son…)". \footcite{2025a}

Les limites de l'IA identifiées, à savoir l'hallucination, la marge d'erreur, les biais culturels, sont ici restreintes puisque les outils d'IA seraient mis à disposition apr skinsoft à des professionnels, et non au grand public directement. Certes, ces outils IA permettront à terme, comme l'indexation automatique, de rendre les collectiosn acessibles au public, mais seulement après vérification d'un professionnel attitré. 

Les professionnels pourraient mettre à profit l'IA pour un double objectif : la préservation et l'accessibilité des collections. \footcite{zotero-item-640}. C'est exactement le cas pour l'indexation automatique. 

\item S-movies 

les collections de film posent des défis spécifiques pour l'implémentation de l'IA, et en même temps, ce serait pour ces collections que l'IA serait pertinente à mettre en place: un défi documentaire 
WEMI 
opacité de la hiérarchisation des niveaux 
refonte de la manière d'indexer dans l'application pour les collections de films : un des enjeux actuels : améliorer l'expérience utilisateur. 

notes depuis la doc interne :
l'IA devra permettre de proposer une indexation pertinente des images par rapport au contenu des métadonnées de chaque fichier et à l'analyse de la représentation de cette image. (reversement des métadonnées). Reconnaissance des types d'objets, extraction des détails visuels, 
Indexer les vidéos ou enregistrements sonores par l'analyse du texte transcrit, indexer les vidéos par l'analyse de l'image, indexer les vidéos ou enregistrements sonores par la reconnaissance vocale, reconnaissance de textes, lien vers les référentiels. 
transcription automatique : convertir les archives orales en texte searchable et els indexer (termes référentiels, autorités, lieux), pointage vers des éléments de référentiels existants (thésaurus). par exemple : citation d'un auteur dans un texte OCR et lien vers la fiche d'auteur. 
Indexation temporelle, ce qui nous intéresse ici : créer des marqueurs automatiques dans les vidéos : changement de scènes, apparition d'objets. séquençage média
Reconnaissance de voix : identifier les intervenants dans les itv
extraction de métadonnées : analyser le contenu sonore (musique, ambiance, silences)
détectin d'incohérence des données : doublons...

\end{itemize}